\section{Methodology}

\textcolor{red}{TO-DO: Clean the folowing paragraph.}

This section begins with the presentation of the Barcelona case study and the description of the datasets employed in the analysis, which serve to illustrate the proposed optimization framework. The framework is then detailed, encompassing the mathematical formulation of the problem through a network-based representation with spatial constraints, followed by a two-step evaluation procedure composed of an MCDM-based scoring stage and a network-level adjustment that balances proximity among stations and overall system accessibility. Next, the topography module is introduced, which further refines distance calculations by accounting for slope effects, and the section concludes with the configuration of the GA used for the optimization.

    \subsection{Data}
    
    This study draws on multiple open data sources available for Barcelona, including the Instituto Nacional de Estadística (INE)~\cite{INE}, OpenStreetMap (OSM)~\cite{OpenStreetMap_web}, Open Data Barcelona (ODB)~\cite{ODB}, and OpenTopoData (OTD)~\cite{OTD}, which provide the following datasets:

    \begin{itemize}

        \item \textbf{Bicycle‐accessible street network:} the bicycle‐accessible street network of Barcelona was obtained by querying OSM via OSMnx~\cite{OSMNx2025} with the network type ‘bike’, yielding a directed graph that comprises approximately 18,700 nodes and 38,000 edges. The graph is formally defined as:
            
            $$ G = (V, E) $$
            
        where $V$ is the set of nodes, with each node $v \in V$ representing street intersections or endpoints; and $E$ is the set of directed edges, where each edge $e \in E$ corresponds to a street link connecting two nodes (Table~\ref{table:notation}). The procedure used to clean and unify this directed graph, removing isolated components and ensuring full bidirectional reachability, is described in Appendix~\ref{app_paper1:graph_cleaning}. The resulting, fully connected bicycle network is shown in Figure~\ref{fig:osm_network}.

        \item \textbf{Locally relevant variables:} alongside the bicycle network, a set of locally relevant variables was collected to characterize each node \(i\in V\). These variables come from the open‐data providers and fall into four categories: socioeconomic (population and income from INE, 2022; education, nationality, household size, vehicle ownership, and unemployment from ODB, 2022–2024), public transport (metro, tram, and bus stops from OSM, accessed 2024), built environment (bike lanes and POIs from OSM, accessed 2024), and topographic (elevation values from OTD, accessed 2024). Table~\ref{table:local_factors_nodes} summarizes each variable’s preprocessing, and full source details appear in Appendix~\ref{app_paper1:data_sources}.  

    \end{itemize}

    
    \begin{figure}[!htbp]
        \centering
        \includegraphics[width=0.65\linewidth]{0_plots/paper1/main/bike_graph.png}
        \caption{\textbf{Barcelona’s OSM bike network after cleaning}. See Appendix~\ref{app_paper1:graph_cleaning} for details on the cleaning procedure.}
        \label{fig:osm_network}
    \end{figure}

    \begin{table}[!htbp]
    \centering
    \caption{\textbf{Notation for the BSS optimization model.}}
    \label{table:notation}
    \footnotesize
    \begin{tabular}{p{0.15\linewidth}p{0.20\linewidth}p{0.60\linewidth}}
        \toprule
        \textbf{Category} & \textbf{Symbol} & \textbf{Definition} \\
        \midrule
        \multirow{3}{*}[0.5em]{\textbf{Graph}} 
        & $G=(V,E)$ & Directed bicycle network with a set of nodes $V$ and a set of edges $E$ \\
        & $d(i,j)$ & Shortest-path distance between nodes $i,j$ \\
        \midrule
        \multirow{5}{*}[0.5em]{\textbf{Stations}} 
        & $S \subset V$ & Set of selected station locations \\
        & $k = |S|$ & Total number of stations \\
        & $x_i \in \{0,1\}$ & Binary decision variable: 1 if node $i$ selected as a station, 0 otherwise \\
        & $X$ & Minimum inter-station distance constraint \\
        & $Y$ & Service coverage radius of the stations \\
        \midrule
        \multirow{4}{*}[0.5em]{\textbf{Local utility}} 
        & $F$ & Total number of local factors \\
        & $V_{\mathrm{norm},f}^i$ & Normalized value of factor $f$ at node $i \in [0,1]$ \\
        & $W_f$ & Weight assigned to factor $f$, with $\sum_{f=1}^F W_f = 1$ \\
        & $U_i$ & Node utility:  \\
        \midrule
        \multirow{4}{*}[0.5em]{\textbf{Network metrics}} 
        & $S_{\mathrm{prox}}$ and $S_{\mathrm{acc}}$ & Inverted normalized proximity and accessibility scores, respectively \\
        & $\alpha \in [0,1]$ & Trade-off parameter balancing proximity and accessibility \\
        & $CS$ & Combined network score: $CS = \alpha \cdot S_{\mathrm{prox}} + (1-\alpha) \cdot S_{\mathrm{acc}}$ \\
        \midrule
        \textbf{Objective} & $U^*_{\mathrm{BSS}}$ & Total BSS utility: $U^*_{\mathrm{BSS}} = \left(\sum_{i \in S} U_i\right) \times CS$ \\
        \bottomrule
    \end{tabular}
    \end{table}

    \subsection{Optimization problem}
    
    Building on the data described above, this study formulates BSS planning as an optimization problem in which station locations are chosen to maximize the system’s total utility (Eq.~\ref{eq:full_formula}). The framework is flexible in that it can incorporate a wide range of spatial factors, adjust weighting schemes, and prioritize proximity among stations or accessibility of the system to reflect different planning priorities.

    The aim of the optimization is to select a set $ S \subset V $ of $ k $ nodes from the OSM bike graph to serve as BSS stations. To do this, each node $i \in V$ is associated with a binary decision variable $x_i$, where $x_i = 1$ if node $i$ is selected as station and $x_i = 0$ otherwise. The only constraint in this optimization is spatial (Eq.~\ref{eq:spatial_constraint}): any two selected stations $ i $ and $ j $ must be separated by at least $X$ meters apart considering the paths of the directed bike network: 

        \begin{equation}\label{eq:spatial_constraint}
            d(i,j) \geq X, \quad \text{if } x_i = 1 \text{ and } x_j = 1
        \end{equation}

    where $ d(i,j) $ is the shortest‐path distance. Minimum inter-station separation distances commonly adopted in the BSS literature range from 250\,m to 500\,m \cite{Kabak2018, MateoBabiano2016, Croci2014}; in this study, a fixed value of  $X$ = 300\,m is adopted to illustrate the presented framework, which can accommodate any alternative value of \(X\) as required. Furthermore, the basic path‐length metric \(d(i,j)\) can be substituted with an elevation‐adjusted “equivalent flat distance”  $d_{\mathrm{eq}}(i,j)$, thereby accounting for uphill and downhill trip effects (see Section~\ref{sec:slopes}). 

    Subject to the spatial‐separation constraint above, this work proposes to optimize station locations to maximize overall system utility, defined by two complementary components: the suitability of individual candidate nodes and the structural performance of the network in terms of system compactness and accessibility. These components can sometimes conflict: a candidate location that scores highly on local factors (e.g. locations with high population density or public transport connectivity) may not improve network compactness or accessibility, while a location that contributes to a more compact or accessible network may have lower local suitability. To reconcile these considerations, the framework integrates both perspectives through a two-step approach:

    \begin{enumerate}
        \item \textbf{Station‐level scoring:} assign each node \(i\) a utility \(U_i\) via a weighted sum of normalized local factors chosen based on the system's planning objectives (e.g. population, number of POIs, income, etc.). Each factor is measured within a configurable circular buffer of radius $Y$ around the node, so that the score reflects the surrounding conditions of the potential station location (see Section~\ref{sec:section_nodes_utility}).

        \item \textbf{Network‐level adjustment:} combine the sum of selected node utilities with a composite network metric that balances station proximity and system accessibility according to a user-defined parameter $\alpha$ (see Section~\ref{sec:section_network_utility}).
    \end{enumerate}

    This two‐step procedure yields a utility score for any station set $S$, but does not itself solve the placement problem. Given the combinatorial nature of the placement problem, a GA is employed to identify high-utility BSS station sets (Section \ref{sec:GA}).


    \subsection{Scoring framework}
        
        \subsubsection{Nodes utility} \label{sec:section_nodes_utility}

        The utility score \(U_i\) for each candidate node \(i\) is computed as a spatial weighted sum over \(F\) factors:

            \begin{equation} \label{eq:node_utility}
                U_i \;=\;\sum_{f=1}^F W_f\,V_{\mathrm{norm},f}^i,
            \end{equation}

        where \(V_{\mathrm{norm},f}^i\) represents the normalized value of factor \(f\) within a buffer of radius \(Y\) meters around node \(i\), and \(W_f\) is the user-defined weight for that factor. By adjusting the weights \(W_f\), planners can shift emphasis among competing objectives, while the buffer radius \(Y\) determines the spatial extent of each node’s local context.

        Table~\ref{table:local_factors_nodes} enumerates the local factors considered in this study. These are meant as examples of factors commonly considered in BSS planning rather than a fixed set: any spatial attribute such as land‐use mix, crime rates, ridership, or parking availability can be incorporated by first measuring its value within a specified buffer around each candidate node and then assigning a weight to each spatial attribute. In this study, absolute totals are used (e.g., the number of residents or the count of amenities within the buffer), so that higher underlying values translate directly into higher node scores. Relative indicators, such as per-capita or per-area measures, could also be incorporated, but they express values in relation to a denominator rather than absolute magnitude, which can lead to different interpretations across planning scenarios. Regardless of the choice, all variables are normalized to ensure comparability. Details on the normalization procedures and before/after distributions for these variables are provided in Appendix~\ref{app_paper1:normalization}. Once the normalized scores have been computed, the BSS utility for a given set of stations is determined by the following formula:

            \begin{equation}\label{eq:BSS_utility_no_metrics}
                    U_{BSS} = \max \sum_{i\in V} x_i \cdot U_i
            \end{equation}

        where $V$ is the set of candidate nodes, $x_i$ is the binary decision variable for node $i$ (1 if selected as a station, 0 otherwise), and $U_i$ is the utility score of node $i$.

        \begin{table}[!htbp]
            \centering
            \caption{\textbf{Local factors considered for BSS station placement}. Data sources are detailed in Appendix~\ref{app_paper1:data_sources}.}
            \label{table:local_factors_nodes}
            \scriptsize
            \renewcommand{\arraystretch}{1.2}
            \begin{tabular}{p{0.15\linewidth} p{0.12\linewidth} p{0.73\linewidth}}
                \toprule
                \textbf{Category} & \textbf{Variable} & \textbf{Preprocessing and aggregation within radius \(Y\)} \\
                \midrule
                Socioeconomics & Population 
                    & Section-level counts apportioned to residential building footprints (area-share, see Appendix~\ref{app_paper1:pop_distribution}); building-level counts summed within \(Y\), with partial overlaps weighted by intersecting area.\\
                & Sex 
                    & Male and female counts per section allocated to buildings and aggregated to nodes using the same method as population. \\
                & Age 
                    & Counts in 10-year bins (10–19, …, 70+) per section allocated to buildings and aggregated to nodes using the same method as population. \\
                & Immigration 
                    & Non-Spanish citizen counts per section redistributed to buildings and aggregated to nodes using the same method as population. \\
                & Education 
                    & Counts by attainment level (primary, secondary, tertiary) per section redistributed to buildings and aggregated to nodes using the same method as population. \\
                & Unemployment 
                    & Neighborhood-level unemployment rate applied to section population aged 19–69; resulting counts redistributed to buildings and aggregated to nodes as population. \\
                & Car ownership 
                    & Motorization index (vehicles per 1 000) converted to absolute counts (index × section population / 1 000), redistributed to buildings and aggregated as population. \\
                & Household size 
                    & Average dwelling size (m²) per section aggregated to each node via area-weighted mean over intersecting sections. \\
                & Income 
                    & Average net income per capita per section aggregated to each node via area-weighted mean over intersecting sections. \\
                \midrule
                
                Public transport  & Metro lines 
                    & Number of distinct metro lines whose entrances lie within \(Y\) of each node. \\
                & Tram lines 
                    & Number of distinct tram lines within \(Y\). \\
                & Bus stops 
                    & Number of bus stops (by line) within \(Y\). \\
                \midrule
                
                Built environment & Bike lanes 
                    & Total length of OSM “cycleway” segments within \(Y\). \\
                & POI count 
                    & Number of POIs within \(Y\) considering the categories: healthcare, culture, tourism, recreation, sport, economic and retail, industrial, green, civic, worship, and education, based on a 15-minute city review~\cite{Papadopoulos2023}. Additionally, an overall count of all POIs is also considered.\\
                & POI diversity 
                    & Shannon entropy of POI categories within \(Y\) (zero when no POIs present). \\
                \bottomrule
            \end{tabular}
        \end{table}


        \subsubsection{Network utility} \label{sec:section_network_utility}

        After computing the node's local utility for a candidate set $S$ (Eq.~\ref{eq:BSS_utility_no_metrics}), the BSS score is complemented by an evaluation of the BSS network structure. This evaluation must account for both the proximity among the chosen stations and the coverage of the BSS service area. These two aspects capture different characteristics of the station placement, as a configuration that maximizes station dispersion does not necessarily ensure optimal coverage, and vice versa. A composite metric is constructed by combining the two individual metrics to provide a unified measure of overall network spatial arrangement:

        \begin{itemize}
            \item \textbf{Proximity metric} measures how close the stations in $S$ are, building on the concept of farness~\cite{Sabidussi1966}, which sums the shortest-path distances from one node to all others, and is thus closely related to the reciprocal of the unnormalized closeness centrality. Farness quantifies how far a node is from all other nodes in a network. Here, this idea is adapted to evaluate the spatial distribution of BSS stations by computing the total pairwise farness within the selected stations:

                \begin{equation}
                    Pro = \sum_{{s,t} \subset S} d(s,t)
                \end{equation}

            where $d(s,t)$ represents the shortest-path distance between stations $s$ and $t$, computed using the edges of the directed graph $G$. A larger $Pro$ value indicates that the selected stations are more dispersed across the network.

            \item \textbf{Accessibility metric} measures how well the stations provide accessibility across the entire network. It is based on the $p$-median problem~\cite{Daskin2015, Revelle2008}, a well-known facility location model that seeks to minimize the sum of distances between demand points and their nearest facility. The metric is computed as the total shortest-path distance from each node $v \in V$ to its nearest station $s \in S$:

                \begin{equation}
                    d(v, S) = \min_{s \in S} d(v,s),
                \end{equation}

            Then, the accessibility metric is defined as:
                
                \begin{equation}
                    Acc = \sum_{v \in V} d(v, S)
                \end{equation}

            Lower values of $Acc$ indicate higher accessibility, meaning that, on average, stations are better positioned to minimize travel distance from any node of $G$.
        \end{itemize}

        Because both $Pro$ and $Acc$ depend on network size and station count, they are first min–max scaled to $[0,1]$ for comparability, and then inverted so that higher values denote higher proximity and higher access, as shown in Eq.~\ref{eq:normalization}. The resulting normalized inverted metrics are denoted as $S_{prox}$ and $S_{acc}$.

            \begin{equation} \label{eq:normalization}
                    M_{\mathrm{inv}} = 1 - \frac{M - M_{\min}}{M_{\max}-M_{\min}}.
            \end{equation}

        Since finding the truly minimum or maximum values of the proximity metric is NP-hard~\cite{Hassin1997, Ravi1994} (it requires choosing a subset of nodes to optimize the sum of pairwise shortest‐path distances), its normalization bounds are approximated using heuristics. Similarly, for the accessibility metric, which is based on the \textit{p}-median problem, determining its optimal values is also NP-hard~\cite{Kariv2006, Hakimi1965}, necessitating the use of heuristic approximations for its bounds. The approximate minimum and maximum values are computed as follows (Figure~\ref{fig:graph_metric_bounds}):

        \begin{figure}[!htbp]
            \centering
            \includegraphics[width=0.8\linewidth]{0_plots/paper1/main/Barcelona_graph_metric_bounds.png}
            \caption{\textbf{Normalization bounds for proximity and accessibility metrics on Barcelona’s bike network with \(k=50\) stations}. Stations are shown in red, and the city boundary is outlined in grey.
            \textbf{(A)} Maximum proximity (\(\mathrm{Pro}_{\max}\)), representing the worst‐case proximity of stations.  
            \textbf{(B)} Minimum accessibility (\(\mathrm{Acc}_{\min}\)), representing the best‐case coverage of all nodes.  
            \textbf{(C)} Maximum accessibility (\(\mathrm{Acc}_{\max}\)), representing the worst‐case coverage.  
            Above each subplot, the inverse‐normalized value of the corresponding metric is shown. Note that the proximity lower bound (Eq.~\ref{eq:pro_min}) is agnostic to the city graph and hence omitted.}
            \label{fig:graph_metric_bounds}
        \end{figure}


        \begin{itemize}
            \item \textbf{Minimum proximity distance}: this is the most compact configuration, where each pair of $k$ stations is exactly $X$\,m apart:
                
                \begin{equation}
                    \mathrm{Pro}_{min} =  \binom{k}{2}X
                        \label{eq:pro_min}
                \end{equation}

            \item \textbf{Maximum proximity distance}: this corresponds to the most dispersed BSS, approximated via a \textit{farthest‐first} heuristic. First, the node with the greatest eccentricity is  chosen as the initial station,
                
                \begin{equation}
                    s_{1} = \arg\max_{v\in V}\;\max_{u\in V} d(v,u),
                    \label{eq:seed}
                \end{equation}
            
            ensuring a peripheral starting point. Then, subsequent stations are added one by one by selecting the node whose minimum distance to the existing set is largest:
                \begin{equation}
                    s_i = \arg\max_{v \notin \{s_1,\dots,s_{i-1}\}} \min_{s \in \{s_1,\dots,s_{i-1}\}} d(v,s).
                    \label{eq:farthest}
                \end{equation}
            
            The resulting total pairwise distance gives the upper bound for proximity:
                \begin{equation}
                    \mathrm{Pro}_{max} = \sum_{\{s,t\} \subset S} d(s,t).
                    \label{eq:pro_max}
                \end{equation}

            \item \textbf{Minimum accessibility distance:} this corresponds to  the scenario in which all nodes in \(V\) are as close as possible to a station. It is approximated via \(k\)-means clustering on node coordinates. A station is placed at the node closest to each cluster centroid forming the station set $S_{\mathrm{acc,min}}$. The bound is:

                \begin{equation}
                    \mathrm{Acc}_{min} \approx \sum_{v\in V} \min_{s \in S_{\mathrm{acc,min}}} d(v,s).
                    \label{eq:acc_min}
                \end{equation}


            \item \textbf{Maximum accessibility distance:} this corresponds to the scenario in which all nodes are as far as possible from their nearest station. Using the same definition in Eq.~\ref{eq:acc_min}, the bound is evaluated with $S_{\mathrm{acc,max}}$ instead of $S_{\mathrm{acc,min}}$. In this case, stations are concentrated in a peripheral area of the network. Starting from a peripheral seed node (Eq.~\ref{eq:seed}), subsequent stations are iteratively placed at the closest available nodes while respecting the minimum separation constraint.
        \end{itemize}
            

        To assess both network proximity and accessibility at the same time, the normalized inverted scores of both are combined into a single score defined as:
            \begin{equation}
                CS(\alpha) = \alpha S_{\text{pro}} + (1 - \alpha) S_{\text{acc}}
            \end{equation}
        
        where $\alpha \in [0,1]$ controls the weighting between proximity and accessibility, $S_{\text{pro}}$ is the inverted normalized $Pro$, and $S_{\text{acc}}$ the inverted normalized $Acc$. A value of $\alpha = 0$ prioritizes accessibility, while $\alpha = 1$ focuses solely on station proximity, and intermediate values allow planners to balance these two objectives according to their priorities.


        Based on this, the overall BSS utility is obtained by extending Eq.~\ref{eq:BSS_utility_no_metrics} to combine individual station scores with network performance: 
            \begin{equation} \label{eq:full_formula}
                U_{\text{BSS}}^* = U_{\text{BSS}} \times \left( \alpha \cdot S_{\text{pro}} + (1 - \alpha) \cdot S_{\text{acc}} \right)
            \end{equation}
        where $U_{\text{BSS}}$ is the sum of individual station scores, ensuring that the final utility reflects both geographical coverage and network proximity, as balanced by $\alpha$.

    \subsection{Distance adjustment based on topography}
    \label{sec:slopes}
    
    Beyond the baseline formulation, it is important to consider how topography influences cycling effort. Since shortest-path distances $d(i,j)$ play a central role in the optimization, they should capture not only geometric length but also the effort associated with elevation changes. Uphill segments demand greater exertion, while downhill sections may ease travel but, in the case of steep declines, can also force cyclists to brake to avoid accidents. To represent these effects, the framework allows the use of an elevation-adjusted ‘equivalent flat distance’ $d_{\mathrm{eq}}(i,j)$. In this formulation, uphill segments are penalised and downhill segments rewarded, yielding distances that more closely approximate the perceived effort of cycling, particularly for m-bikes.

    The equivalent flat distance for a single edge \(\ell\) of horizontal length \(d_{\ell}\) and elevation change \(\Delta h_{\ell}\) is computed in the following steps. First, the slope of the edge \(\ell\) is computed and decomposed into uphill and downhill components in one step:
        \[
        g_{+,\ell} = \max\!\bigl(g_{\ell},0\bigr), 
        \quad
        g_{-,\ell} = \min\!\bigl(g_{\ell},0\bigr),
        \quad\text{where}\quad
        g_{\ell} = \frac{\Delta h_{\ell}}{d_{\ell}}.
        \]
    To align with standard cycling guidelines, which recommend a maximum slope of 10\%, \(g_{\ell}\) is then clamped to the interval $[-0.10,0.10]$ \cite{AASHTO1999bicycle, CycleHighways2025, CROW2006}. Applying Parkin \& Rotheram’s linear slope–speed relation \cite{Parkin2010}, the segment speed becomes  
        \[
        v(g_{\ell})=v_{0}+b_{\mathrm{up}}\,g_{+,\ell}+b_{\mathrm{down}}\,g_{-,\ell},
        \]
    with \(v_{0}=6.01\) m/s, \(b_{\mathrm{up}}=-40.02\), and \(b_{\mathrm{down}}=-23.79\) (m/s per unit slope). The resulting travel time \(t_{\ell}=d_{\ell}/v(g_{\ell})\) is then converted back into distance at flat‐ground speed via \(d_{\mathrm{eq},\ell}=t_{\ell}\,v_{0}\). Summing these adjusted lengths over the shortest‐path \(P_{ij}\) between nodes \(i\) and \(j\) yields  
        \[
        d_{\mathrm{eq}}(i,j)=\sum_{\ell\in P_{ij}}d_{\mathrm{eq},\ell},
        \]
    which may replace the original \(d(i,j)\) in the spatial constraint whenever slope effects are deemed relevant. 

    Thus, the procedure provides an adjusted distance matrix where each shortest path $d_{eq}(i,j)$ reflects slope effects. This measure can directly replace the original distances $d(i,j)$ in the spatial constraint, penalizing uphill segments and rewarding downhill ones.

    \subsection{Optimization modelling}
    \label{sec:GA}

    In this study, a GA is employed to optimize the placement of BSS stations. GAs are population-based search procedures inspired by natural selection and genetics, well suited to exploring large combinatorial spaces and handling multiple objectives and constraints without requiring differentiable or convex formulations. These properties make them appropriate for station location problems, where the search space is discrete, constraints such as minimum station spacing must be enforced, and solutions must adapt to different planning objectives while retaining the same underlying spatial structure. A GA evolves a population of fixed-length ‘chromosome’ candidates, each encoding all decision variables, through the operators of \textit{selection}, which favors higher-fitness individuals; \textit{crossover}, which splices gene segments from two parents to create offspring; and \textit{mutation}, which randomly alters individual genes to maintain diversity and avoid premature convergence. This evolutionary cycle continues until a stopping criterion such as a maximum number of generations or a convergence threshold is met. For a comprehensive overview of GA concepts and variants, the reader is referred to Katoch et al.~\cite{Katoch2021}.

    Our algorithm follows a single-objective, non-adaptive GA with elitism. Each chromosome is a binary vector over all candidate sites (‘station’ versus ‘no station’), encoding a complete BSS deployment plan $S$ with $k$ stations. The initial population of $N$ candidate plans is generated by constrained random sampling: a first station is selected at random, and subsequent stations are added only if they lie at or beyond the prescribed minimum distance $X$ from all previously chosen sites, until $k$ stations are placed. The population then evolves through selection, crossover, and mutation, while elitism ensures that a fixed fraction of the best solutions is preserved across generations. 

    This design was chosen for its robustness and compatibility with custom, constraint-aware operators. Although other variants such as steady-state, adaptive, or multi-objective GA variants exist, the chosen formulation proved sufficient to capture the spatial structure of the problem and to generalize across diverse planning scenarios. As demonstrated in Section~\ref{GA_vs_MILP}, this consistency allows a single hyper-parameter optimization process to be reused across all experimental configurations.


        \subsubsection{Hyper-parameter tuning}

        To ensure adaptability across diverse planning objectives, a comprehensive hyper-parameter search was conducted on a representative scenario (\(k=60\) with a randomly generated local-factor weight vector), hypothesizing that these hyper-parameter settings would generalize to other station counts and weight combinations (see Section~\ref{GA_vs_MILP}). Each hyper-parameter configuration was assessed in terms of solution quality and population diversity over time. The main components of the algorithm, along with the alternative configurations explored, are detailed below and in Table~\ref{table:ga_grid}. 

        {%
        \setlength{\parskip}{0.45\parskip}%
        \begin{itemize}
        \item \textbf{Population size.} The algorithm starts with $N$ solutions, each comprising $k$ station locations that satisfy $X$. A larger $N$ enhances diversity, giving the algorithm a broader view of the solution space and increasing the chances of escaping local optima, but it also increases computational time.

        \item \textbf{Selection strategy.} Determines how parent solutions are selected for reproduction:
        
        \begin{itemize}
            \item \textit{Roulette-wheel:} individuals are selected with a probability proportional to their fitness score. High-performing solutions are more likely to be chosen, but lower-quality ones retain a non-zero chance.
            \item \textit{Tournament:} a small subset of individuals is randomly drawn from the population, and the one with the highest fitness within this group is selected.  This approach introduces selection pressure while still allowing less-fit individuals a chance to be chosen.            
        \end{itemize}

        \item \textbf{Crossover strategy.} Controls how offspring are generated from parent solutions: 
        
        \begin{itemize}
            \item \textit{Greedy:} candidate nodes from both parents are ranked by score, and the best ones are iteratively added while respecting spatial constraints. 
            \item \textit{Top-first:} the top 25\% of nodes by score are shuffled and selected first, with remaining slots filled by random picks from the other 75\%.
            \item \textit{Weighted-random:} candidates are sampled with probabilities proportional to their scores, balancing exploration and exploitation. 
        \end{itemize}

        \item \textbf{Mutation strategy.} each node of non-elite offspring has a probability of mutation, capped at 20\% of nodes. Node replacements must satisfy the distance constraint \(X\).

        \item \textbf{Elitism.} A fixed proportion of the best-performing solutions is directly transferred to the next generation. This guarantees that high-quality individuals are preserved and not lost through random variation.

        \item \textbf{Early stopping.} To avoid unnecessary computation, the algorithm is terminated if no improvement is observed in either the best or average fitness score after 300 consecutive generations.
        \end{itemize}
        }%


        \begin{table}[!htbp]
        \centering
        \caption{\textbf{GA grid search hyper-parameter values}. Combinations of hyperparameters that were not feasible due to insufficient numbers in the different GA operators were skipped.}
        \label{table:ga_grid}
        \footnotesize
            \begin{tabular}{ll}
            \toprule
            \textbf{Hyperparameter} & \textbf{Values tested} \\
            \midrule
            Population size & 25, 50, 100, 200 \\
            Mutation rate & 0.01, 0.02, 0.05, 0.1, 0.2 \\
            Elite fraction & 0.01, 0.02, 0.05, 0.1, 0.2 \\
            Selection strategy & Tournament, Roulette-wheel \\
            Crossover strategy & Greedy, Weighted-random, Top-first \\
            Max generations & 10,000 (with early stopping)\\
            \bottomrule
            \end{tabular}
        \end{table}

        The best‐performing set comprised tournament selection, greedy crossover, \(N=50\), mutation rate 0.05, and elite fraction 0.05.  


        \subsubsection{GA accuracy bench-marking}
        \label{GA_vs_MILP}

        Although the GA hyper-parameters were tuned on a single scenario, they are expected to generalize across planning scenarios, as the network topology remains fixed. To validate this hypothesis, GA solution utilities have been compared against the exact optima obtained from a Mixed Integer Linear Programming (MILP) model. GAs are meta-heuristic and cannot guarantee global optimality, whereas MILP solvers exhaustively search the binary decision space for linear objectives and certify the optimal solution, making them ideal benchmarks.

        Because the full BSS utility combines local-factor and network-structure metrics in a non-linear objective, it cannot be solved exactly by MILP. Instead, the MILP benchmark optimizes only the linear component
        \[
        f_{\mathrm{MILP}} = \max_{x}\;\sum_{i\in V} x_i\,U_i
        \]
        subject to  
        \[
        \sum_{i\in V} x_i = k,\quad
        x_i + x_j \le 1\;\;\forall\,(i,j)\in\mathcal{E},\quad
        x_i\in\{0,1\},
        \]
        where \(V\) is the set of candidate nodes, \(U_i\) the local-factor utility of node \(i\), \(k\) the prescribed number of stations, and \(\mathcal{E}=\{(i,j)\mid d(i,j)<X\}\) enforces the minimum inter-station distance \(X\). This linear formulation preserves the same stations' deployment space while substantially reducing computational complexity and avoiding approximation errors required to linearize non-linear terms. It should be emphasized, however, that this MILP benchmark is based on a reduced formulation of the problem that excludes the non-linear component (i.e., incorporating the network-structure metrics into $U_{\mathrm{BSS}}$). As such, it provides only a partial measure of GA optimality. To enable a fair comparison, GA runs were also applied to obtain $U_{\mathrm{BSS}}$. GA robustness was assessed using two experiments: 

        \begin{enumerate}
        \item \textbf{Weight sensitivity} (fixed $k=60$).
            A total of 118 weight vectors for the local factors were generated by Monte Carlo sampling. A sample size of 118 was chosen based on a preliminary power analysis with 20 samples to ensure that a one-sided test at the 0.5\% gap threshold would achieve at least 90\% power with \(\alpha=0.05\). For each weight vector, the MILP solved the linear model to optimality and the GA was executed with the tuned hyper-parameters, enabling paired comparisons of \(U_{\mathrm{BSS,GA} }\) and \(U_{\mathrm{BSS,MILP}}\).

        \item \textbf{Station-count sensitivity (fixed weights, 24 vectors).}  
            Four station counts (\(k=30,60,90,120\)) were each tested against 96 randomly generated weight vectors. The choice of 96 weight combinations per \(k\) is guided by the same power analysis.

        \end{enumerate}

        For every test instance, the optimality gap  

        \[
        \mathrm{Gap} = 100\times\frac{f_{\mathrm{BSS,MILP}} - f_{\mathrm{BSS,GA}}}{f_{\mathrm{BSS,MILP}}}
        \]

        was computed as in Alkhalifa et al.~\cite{Alkhalifa2025} and runs with $\mathrm{Gap}\leq1\%$ were classified as near-optimal, following the convention in the heuristic optimization literature \cite{Lagos2023, Madadi2024, Pisinger2007}. Results are summarized in Table~\ref{table:MILP_combined}. In all instances the GA achieved gaps under 1\% of the MILP optimum, confirming that neither weight variation nor station count degrades GA accuracy under the simplified linear formulation. To statistically validate this equivalence, paired percentage gaps were tested for normality (Shapiro--Wilk) and, depending on the result, subjected to either one-sample $t$-tests or Wilcoxon signed-rank tests against a $0.5\%$ threshold. In every scenario the one-tailed $p$-value was below $0.001$, and observed mean gaps ranged from $0.23\%$ to $0.35\%$. 
        %Post-hoc power analyses ($\alpha=0.001$, power=0.9999) showed that our sample sizes far exceed the numbers needed to detect these effects. 

        These findings demonstrate that the GA can approximate exact solutions very closely when the problem is linear. Since the MILP benchmark excludes the network metric of the full BSS utility $U^*_{\mathrm{BSS}}$, the results should be interpreted as an indicative rather than exhaustive validation. On the other hand, although GA runtimes (10–66 min) exceed MILP times (6 min), only the GA approach, among the two, can optimize the full $U^*_{\mathrm{BSS}}$, and its near-optimal performance on the linear proxy gives confidence in its ability to tackle the complete problem.


        \begin{table}[!htbp]
        \centering
        \caption{\textbf{GA vs.\ MILP benchmark results.} Results for both the weight and station count sensitivity analysis. GA scores correspond to the best solution of the population.}
        \label{table:MILP_combined}
        \footnotesize
        \begin{tabular}{r r r cc cc}
        \toprule
        & & & \multicolumn{2}{c}{\textbf{MILP}} & \multicolumn{2}{c}{\textbf{GA}} \\
        \cmidrule(lr){4-5} \cmidrule(lr){6-7}
        \textbf{Scenario} & \textbf{\(k\)} & \textbf{Weights} 
            & \textbf{Score (Avg\,$\pm$\,Std)} & \textbf{Time (Avg\,$\pm$\,Std)} 
            & \textbf{Score (Avg\,$\pm$\,Std)} & \textbf{Time (Avg\,$\pm$\,Std)} \\
        \midrule
        Weights & 60 & 118 & 53.48 $\pm$ 0.51 & 5.92 $\pm$ 0.07 & 53.29 $\pm$ 0.54 & 33.60 $\pm$ 10.25 \\
        \midrule
        Stations & 30 & 24  & 27.73 $\pm$ 0.10 & 5.92 $\pm$ 0.08 & 27.67 $\pm$ 0.13 & 10.51 $\pm$ 2.76 \\
                            & 60 & 24  & 53.44 $\pm$ 0.43 & 5.93 $\pm$ 0.11 & 53.26 $\pm$ 0.46 & 21.45 $\pm$ 4.69 \\
                            & 90 & 24  & 77.77 $\pm$ 0.77 & 5.94 $\pm$ 0.06 & 77.50 $\pm$ 0.76 & 38.01 $\pm$ 11.86 \\
                            & 120& 24  & 100.89 $\pm$ 1.03 & 5.96 $\pm$ 0.10 & 100.60 $\pm$ 1.03 & 65.66 $\pm$ 20.23 \\
        \bottomrule
        \end{tabular}
        \end{table}